\section{双用途风险：Agent 能力提升如何改变网络安全攻防}

课程最后转向 cyber security dual-use 问题：AI 可以同时增强攻击者和防守者，关键是 ``谁受益更快''。讲者展示的研究结论较为审慎：短期内由于攻击与防守的自然不对称，前沿 AI 往往更容易先提升攻击效率，因此需要持续监测能力变化而不是一次性判断。

\begin{importantbox}{为什么必须做 Continuous Monitoring}
Agent 能力曲线变化快，静态 benchmark 很快过时。持续评测能帮助社区及时发现：
\begin{itemize}
    \item 新能力出现是否导致攻击成功率跃迁；
    \item 防御策略是否在新模型下失效；
    \item 哪些任务上 ``攻击增益'' 快于 ``防守增益''。
\end{itemize}
\end{importantbox}

讲者提到的 cyber benchmark（大规模任务集合）和 observatory 机制，实质上是把 ``一次论文评测'' 升级成 ``持续治理基础设施''。这对政策和产业都有现实意义：只有连续观测，才能决定何时需要更强默认防护、何时需要限制高风险能力暴露。

\subsection{本章小结}
双用途风险决定了 Agent 安全不是纯技术闭环，还涉及生态协作与治理节奏。持续监测比一次性结论更符合现实。

